% GENERATED FILE. Edit canonical lessons, then run npm run generate:books.
% canonical-source-hash: fnv1a64:5406942ec2e0c2a0
% canonical-lessons: IT-C298-banda, IT-C298-complice, IT-C298-aggredire, IT-C298-aggressione, IT-C298-minacciare

\chapter{La banda}
\label{ch:it-la-banda}
\hlchaptermodality{298}

\begin{chapteropening}
\textbf{By the end of this chapter:} \emph{I can talk about a gang and its accomplice, an attack and a threat.}

Complete the last lesson of chapter 298: I can talk about a gang and its accomplice, an attack and a threat.
\end{chapteropening}

\section[la banda]{la banda — the gang; the brass band}

\label{lesson:IT-C298-banda}



\begin{quote}
Before the new one: say the Italian for to capture, then the Italian for to flee.
\end{quote}

\subsection*{You'll want to know: la banda}
\textbf{la banda} — \textquotedblleft{}the gang; the brass band\textquotedblright{}.

\begin{quote}
Una banda ha rapinato tre negozi.
\end{quote}

\begin{quote}
\emph{A gang robbed three shops.}
\end{quote}

At a town festival, \emph{la banda} is the brass band.

\begin{cousinweb}
\emph{banda} \textquotedblleft{}a band of people\textquotedblright{} is usually traced to Provençal \emph{banda} \textquotedblleft{}a troop\textquotedblright{}, which EtymDB derives from a Frankish word \emph{binda} \textquotedblleft{}a bond, a link\textquotedblright{}.
\end{cousinweb}

\begin{grammarlens}[title={What you've built}]
One more word for telling, explaining and describing at B1.
\end{grammarlens}

\subsection*{Your turn}
\begin{itemize}
\raggedright
  \item \emph{Say it:} \emph{la banda}
  \item \emph{Say it:} \emph{la banda}, once more
  \item \emph{Say it:} a sentence of your own with \emph{la banda}
  \item \emph{Recall:} say the Italian for to capture, then the Italian for to flee, then say \emph{la banda} again
\end{itemize}

\begin{tcolorbox}[breakable,colback=teal!4,colframe=teal!35!black,title={Before you move on}]
What does la banda mean? (\textquotedblleft{}The gang; the brass band\textquotedblright{}.) Say it once more.
\end{tcolorbox}
\section[il complice]{il complice — the accomplice}

\label{lesson:IT-C298-complice}



\begin{quote}
Before the new one: say the Italian for to flee, then the Italian for the gang.
\end{quote}

\subsection*{You'll want to know: il complice}
\textbf{il complice} — \textquotedblleft{}the accomplice\textquotedblright{}.

\begin{quote}
Il rapinatore aveva un complice in macchina.
\end{quote}

\begin{quote}
\emph{The robber had an accomplice in the car.}
\end{quote}

Also an adjective: \emph{uno sguardo complice}, a knowing look.

\begin{cousinweb}
\emph{complice} is a learned borrowing from Late Latin \emph{complex, complicis} \textquotedblleft{}closely joined, an accomplice\textquotedblright{}, from \emph{complicāre} \textquotedblleft{}to fold together\textquotedblright{} (Lewis \& Short).
\end{cousinweb}

\begin{grammarlens}[title={What you've built}]
One more word for telling, explaining and describing at B1.
\end{grammarlens}

\subsection*{Your turn}
\begin{itemize}
\raggedright
  \item \emph{Say it:} \emph{il complice}
  \item \emph{Say it:} \emph{il complice}, once more
  \item \emph{Say it:} a sentence of your own with \emph{il complice}
  \item \emph{Recall:} say the Italian for to flee, then the Italian for the gang, then say \emph{il complice} again
\end{itemize}

\begin{tcolorbox}[breakable,colback=teal!4,colframe=teal!35!black,title={Before you move on}]
What does il complice mean? (\textquotedblleft{}The accomplice\textquotedblright{}.) Say it once more.
\end{tcolorbox}
\section[aggredire]{aggredire — to attack, to assault}

\label{lesson:IT-C298-aggredire}



\begin{quote}
Before the new one: say the Italian for the gang, then the Italian for the accomplice.
\end{quote}

\subsection*{You'll want to know: aggredire}
\textbf{aggredire} — \textquotedblleft{}to attack, to assault\textquotedblright{}.

\begin{quote}
Due uomini lo hanno aggredito di notte.
\end{quote}

\begin{quote}
\emph{Two men attacked him at night.}
\end{quote}

Present like \emph{capire}: \emph{aggredisce}.

\begin{cousinweb}
\emph{aggredire} is a learned borrowing from Latin \emph{aggredī} \textquotedblleft{}to go at, to attack\textquotedblright{} (EtymDB), from \emph{ad} and \emph{gradī} \textquotedblleft{}to step\textquotedblright{}.
\end{cousinweb}

\begin{grammarlens}[title={What you've built}]
One more word for telling, explaining and describing at B1.
\end{grammarlens}

\subsection*{Your turn}
\begin{itemize}
\raggedright
  \item \emph{Say it:} \emph{aggredire}
  \item \emph{Say it:} \emph{aggredire}, once more
  \item \emph{Say it:} a sentence of your own with \emph{aggredire}
  \item \emph{Recall:} say the Italian for the gang, then the Italian for the accomplice, then say \emph{aggredire} again
\end{itemize}

\begin{tcolorbox}[breakable,colback=teal!4,colframe=teal!35!black,title={Before you move on}]
What does aggredire mean? (\textquotedblleft{}To attack, to assault\textquotedblright{}.) Say it once more.
\end{tcolorbox}
\section[l'aggressione]{l'aggressione — the attack, the assault}

\label{lesson:IT-C298-aggressione}



\begin{quote}
Before the new one: say the Italian for the accomplice, then the Italian for to attack.
\end{quote}

\subsection*{You'll want to know: l'aggressione}
\textbf{l'aggressione} — \textquotedblleft{}the attack, the assault\textquotedblright{}.

\begin{quote}
L'aggressione è avvenuta davanti a un bar.
\end{quote}

\begin{quote}
\emph{The attack took place outside a bar.}
\end{quote}

A news word.

\begin{cousinweb}
\emph{aggressione} is a learned borrowing from Latin \emph{aggressiō} \textquotedblleft{}a going towards, an onset\textquotedblright{}, from \emph{aggredī} \textquotedblleft{}to go at, to attack\textquotedblright{} (Lewis \& Short).
\end{cousinweb}

\begin{grammarlens}[title={What you've built}]
One more word for telling, explaining and describing at B1.
\end{grammarlens}

\subsection*{Your turn}
\begin{itemize}
\raggedright
  \item \emph{Say it:} \emph{l'aggressione}
  \item \emph{Say it:} \emph{l'aggressione}, once more
  \item \emph{Say it:} a sentence of your own with \emph{l'aggressione}
  \item \emph{Recall:} say the Italian for the accomplice, then the Italian for to attack, then say \emph{l'aggressione} again
\end{itemize}

\begin{tcolorbox}[breakable,colback=teal!4,colframe=teal!35!black,title={Before you move on}]
What does l'aggressione mean? (\textquotedblleft{}The attack, the assault\textquotedblright{}.) Say it once more.
\end{tcolorbox}
\section[minacciare]{minacciare — to threaten}

\label{lesson:IT-C298-minacciare}



\begin{quote}
Before the new one: say the Italian for to attack, then the Italian for the attack.
\end{quote}

\subsection*{You'll want to know: minacciare}
\textbf{minacciare} — \textquotedblleft{}to threaten\textquotedblright{}.

\begin{quote}
Lo hanno minacciato con un coltello.
\end{quote}

\begin{quote}
\emph{They threatened him with a knife.}
\end{quote}

\emph{Minacciare di} + infinitive: \emph{minaccia di piovere}, it looks like rain.

\begin{cousinweb}
\emph{minacciare} is built on \emph{minaccia} \textquotedblleft{}threat\textquotedblright{}, from Latin \emph{mināciae} \textquotedblleft{}threats\textquotedblright{}, built on \emph{mināx} \textquotedblleft{}threatening\textquotedblright{} (Lewis \& Short).
\end{cousinweb}

\begin{grammarlens}[title={What you've built}]
One more word for telling, explaining and describing at B1.
\end{grammarlens}

\subsection*{Your turn}
\begin{itemize}
\raggedright
  \item \emph{Say it:} \emph{minacciare}
  \item \emph{Say it:} \emph{minacciare}, once more
  \item \emph{Say it:} a sentence of your own with \emph{minacciare}
  \item \emph{Recall:} say the Italian for to attack, then the Italian for the attack, then say \emph{minacciare} again
\end{itemize}

\begin{tcolorbox}[breakable,colback=teal!4,colframe=teal!35!black,title={Before you move on}]
What does minacciare mean? (\textquotedblleft{}To threaten\textquotedblright{}.) Say it once more.
\end{tcolorbox}
